{\bf [40 points] Clustering}\\

In this question, we will perform and compare two clustering algorithms on real single-cell RNA-seq data: (i) K-means clustering and (ii) graph-based Leiden clustering. The goal is to understand how different clustering assumptions affect both quantitative performance and biological interpretability in single-cell analysis.

\textbf{Data Description}

We provide a single-cell RNA-seq dataset in an AnnData object (file: \texttt{sc\_pca.h5ad}). The data are derived from:

\begin{quote}
Smillie, Christopher S., et al. \emph{``Intra- and inter-cellular rewiring of the human colon during ulcerative colitis.''} Cell 178.3 (2019): 714--730.
\end{quote}

The AnnData object contains:

\begin{itemize}
    \item \textbf{Cells:} 34{,}162 cells.
    \item \texttt{obsm} (embeddings): \texttt{X\_pca}, a 30-dimensional PCA embedding for each cell.
    \item \texttt{obs} (cell metadata): a column \texttt{merged\_annotation} with 9 cell type labels:
    Plasma, Follicular, CD4+ Memory, CD8+ LP, Enterocytes, TA, Macrophages, Immature Goblet, Immature Enterocytes 1.
\end{itemize}

All clustering should be performed using the PCA embeddings (\texttt{X\_pca}) only. The cell type annotations are provided \emph{for evaluation and interpretation purposes only}.

\textbf{Notes:}
\begin{itemize}
    \item You may use \emph{any} programming language and packages (e.g., Scanpy, Seurat, scikit-learn) that you find useful.
    \item No skeleton code is provided. Your code will not be graded by an autograder; instead, you should submit your code together with your written solutions. Partial credit may be awarded for some correct implementations even if final results are incorrect.
    \item For K-means clustering, you should use K-means++ initialization or multiple random restarts and select the best solution.
\end{itemize}

\vspace{20pt}

\textbf{Your task}

\begin{enumerate}

\item (15 points) \textbf{K-means clustering.}

K-means clustering partitions the data into $k$ clusters by minimizing the within-cluster sum of squared errors (SSE) in Euclidean space. Given cluster assignments $\{C_1, \dots, C_k\}$ and cluster centroids $\{\mu_1, \dots, \mu_k\}$, the SSE is defined as:
\[
\mathrm{SSE} = \sum_{c=1}^{k} \sum_{x_i \in C_c} \lVert x_i - \mu_c \rVert^2 .
\]

\begin{enumerate}
    \item (5 points) Run K-means clustering for $k = 1,2,\dots,15$ on \texttt{X\_pca}. Plot SSE versus $k$ and identify the elbow point. Briefly explain what the elbow point represents and why it may be ambiguous for single-cell data.
    
    \item (5 points) Run K-means with $k=10$. To ensure stable results, use either K-means++ initialization or run K-means multiple times with different random initializations and select the solution with the lowest SSE. Report the final SSE value.
    
    \item (5 points) Compute the Adjusted Rand Index (ARI) between the K-means cluster assignments ($k=10$) and the ground-truth labels \texttt{merged\_annotation}. The ARI between two partitions $U$ and $V$ is defined as:
    \[
    \mathrm{ARI} = \frac{\mathrm{RI} - \mathbb{E}[\mathrm{RI}]}{\max(\mathrm{RI}) - \mathbb{E}[\mathrm{RI}]},
    \]
    where RI is the Rand Index measuring pairwise agreement between clusterings. Report and interpret the ARI value.
\end{enumerate}

%%%%%%%%%%%%%%%%%%
%%%%%%%%%%%%%%%%%%

\vspace{20pt}

\item (15 points) \textbf{Graph-based Leiden clustering.}

Graph-based clustering methods operate on a neighborhood graph and can capture clusters with arbitrary shapes. In this part, you will apply Leiden clustering.

\begin{enumerate}
    \item (5 points) Construct a $k$-nearest neighbor (kNN) graph from \texttt{X\_pca} using $k=20$ neighbors per cell. Clearly state the distance metric used.
    
    \item (5 points) Run the Leiden algorithm on the kNN graph with resolution parameter \texttt{resolution = 0.1}. Report the number of clusters obtained.
    
    \item (5 points) Compute the ARI between the Leiden cluster assignments (resolution $=0.1$) and \texttt{merged\_annotation}. Report and interpret the ARI value.
\end{enumerate}

%%%%%%%%%%%%%%%%%%
\begin{solution}

\end{solution}
%%%%%%%%%%%%%%%%%%

\vspace{20pt}

\item (10 points) \textbf{Visualization and comparison.}

UMAP is commonly used to visualize single-cell data and clustering results.

\begin{enumerate}
    \item (5 points) Plot a UMAP embedding colored by K-means clusters using $k=4$.
    \item (5 points) Plot the same UMAP embedding colored by Leiden clusters using \texttt{resolution = 0.01}.
\end{enumerate}

\textbf{Question:} Based on the UMAP visualizations of K-means ($k=4$) and Leiden (resolution $=0.01$), what is a potential disadvantage of K-means–based clustering methods in single-cell RNA-seq analysis? Briefly explain your reasoning.

%%%%%%%%%%%%%%%%%%
\begin{solution}

\end{solution}
%%%%%%%%%%%%%%%%%%

\end{enumerate}

\vspace{10pt}


